The vortex lattice method was in name introduced by \textsc{Falkner} in a series of papers during the tailend of the 1940s.%
\footnote{%
        \cite{falkner1946accuracy} \textsc{Falkner} (1946)
}%
\textsuperscript{,}%
\footnote{%
        \cite{falkner1953solution} \textsc{Falkner} (1947)
}%
\textsuperscript{,}%
\footnote{%
        \cite{falkner1952scope} \textsc{Falkner} (1949)
}
%
This formulation, however, is entirely unsuited for unsteady vortex systems, as he therein postulates the vortex system to assume the form of \textsc{Prandtl}'s steady horseshoe vortex system.
It is formulated as follows.
Consider a wing, be it swept or not, and section it spanwise with respect to the unswept axis into segments of ideally equal width.
Place so, according to constant proportion of the chord, lines from root to tip of the wing at one quarter and three quarters of the chord, measured from leading edge to trailing edge, orthogonally to the segmentation of the wing.
Draw then a vortex filament between the edges of the segments, through the quarterchord line for each segment.
Emanating therefrom, along the edges of the segments towards the trailing edge and beyond to infinity, are vortex filaments.
At infinity, we consider filaments meeting each of these emanating filaments, forming a closed vortex loop for each segment.
The collection of these loops is what is known as the vortex lattice.
At the threequarterschord, one will find the induced velocity of the vortex system, enforcing impermeability.
For illustrative purposes, the unimaginative reader is referred to figure 1 in \textsc{Falkner}'s 1947 report.

The reader would be remiss to not consider this business of extending the vortex filaments along the edges to infinity initially suspect in a discussion on unsteady aerodynamic schemes.
For it is indeed the case that this initial formulation of the vortex lattice method is inherently steady.
We have already discussed the fact that the inviscid formulation of the interaction between fluid flow and some solid object cannot account for the generation of vortical motion.
As an \emph{ad hoc} heuristic, the matter is often explained away by the prior shedding of a vortex, now carried downstream to infinity.
Thus by the circulation theorem \ref{thm:circulation-theorem}, the vorticity in the vicinity of the solid object must be of equal but negative magnitude, so that the total circulation of the system remains constant.
Whence the vortex filament at infinity.
Confer with the explanation of \textsc{Newman}.%
\footnote{%
        \cite{newman2017marine} \textsc{Newman}, \S5.1, pp.169--172
}

This peculiar placement of the bound vortex filament and control point for the evaluation of the downwash was first introduced in an article by \textsc{Pistolesi},%
\footnote{%
        \cite{deyoung1976historical} \textsc{DeYoung} (1976)
}
which is unfortunately completely unavailable to the author.
These placements were later shown to be optimal, even for wings of multiple chordwise panels, the vortex being placed at the quarterpoint, and the collocation point at the threequarterspoint.%
\footnote{%
        \cite{scholz1950beitrage} \textsc{Scholz} (1950)
}%
\textsuperscript{,}%
\footnote{%
        \cite{byrd1951erganzung} \textsc{Byrd} (1951)
}


An albeit naïve yet illustrative explanation of the modern vortex lattice method is the exchange of the doublets of the standard boundary element method for vortex filaments along their boundaries.
The mathematical foundation for this apparent equivalence between a dipole panel and a vortex panel is elusive in the literature.
This relationship was to the author's knowledge first demonstrated by \textsc{Hess},%
\footnote{%
        \cite{hess1972calculation} \textsc{Hess} (1972), Appendix A
}
\textsc{Morino} proving the very same in a later work,%
\footnote{%
        \cite{morino1993boundary} \textsc{Morino} (1993), Appendix C
}
relating the tangential derivative of the jump discontinuity in the potential over the boundary $q$ to vorticity on the boundary.
The intrinsic object on such a surface is indeed the exterior derivative $\dee q$, its \textsc{Hodge} dual accordingly aligning with the previously discussed notion of vortex sheets.
When the wake and bound panels have piecewise constant values for the potential jump, this exterior derivative vanishes in the classical sense.
The weak derivative, however, admits us to study the relationship between this potential jump and data at the edges of the panel through the fundamental theorem of exterior calculus.%
\footnote{%
  \cite{bonaldi2025exterior} \textsc{Bonaldi} \emph{et al.} (2025)
}
Admittedly, the jumps in potential are treated in greater detail through the language of currents in the framework finite element exterior calculus,%
\footnote{%
  \cite{licht2017complexes} \textsc{Licht} (2017)
}\textsuperscript{,}%
\footnote{%
  \cite{chen2025distributional} \textsc{Chen}, \textsc{Huang}, \& \textsc{Zhang} (2025)
}
which we shall lend credence to, but not discuss in any meaningful detail.
Nevertheless, the weak derivative exists in a weak sense on the boundary skeleton.
For the wake, we test the \textsc{Hodge} dual of the panelwise constant jump in potential against a smooth one-form, yielding

\begin{equation}
  {(\afrak, \star\dee q)}_{S_{\mathrm{w}}} \equiv - \sum_{\mathrm{p}} q_{\mathrm{p}} \int_{\mathrm{p}} \dee\afrak = - \sum_{\mathrm{p}} q_{\mathrm{p}} \int_{\partial\mathrm{p}} \afrak = \sum_{\mathrm{e}} \Gamma_{\mathrm{e}} \int_{C_{\mathrm{e}}} \afrak(\tvec_{\mathrm{e}}) \,\dee s.
\end{equation}

\noindent
We have here used the fundamental theorem of multivariate calculus, and our preliminary work on vortical fluid flows, after first defining the test over the portion of the boundary that is the wake in a panelwise manner.
Interior edges must thus be collected twice, with opposite orientations, yielding a difference on the edge that is exactly equivalent to the edge circulation.
Edges whose neighboring panel does not exist, clearly yields no difference in potential jump from the panel value.
In explicit mathematical terms, we have that for two adjacent panels $\mathrm{p}_1$ and $\mathrm{p}_2$, sharing the edge $\mathrm{e}$ between them, their respective values for the jump in potential over their interface being $q_{\mathrm{p}_1}$ and $q_{\mathrm{p}_2}$, we define up to orientation the circulation on the shared edge $\Gamma_{\mathrm{e}} = q_{\mathrm{p}_2} - q_{\mathrm{p}_1}$, in a weak sense.
Henceforth, we shall consider the induced velocity from a single panel

\[
  \uvec(\bzhe) \approx \frac{1}{4\pi} \sum_{i = 1}^{4} \Gamma_i \int_{C_i} \frac{\xvec - \bzhe}{r^3} \times \dee\svec,
\]

\noindent
where the circulation $\Gamma_i = q_{\mathrm{p}}$.
For each filament, it is possible calculate the induced velocity analytically.
In order to reduce clutter, we drop the index indicating filament.
The procedure of evaluating for one filament in a lattice is illustrated in figure \ref{fig:line_integral}.

\begin{figure}[h]
        \centering
        \scalebox{1}{\includegraphics{line_integral.pdf}}
        \caption{The calculation of the induced velocity from the vortex filament parametrized by $\xvec(s)$.}
        \label{fig:line_integral}
\end{figure}

\noindent
We consider the parametrization of the vortex filament with respect to the midpoint $\bkha$, with endpoints $\xvec^\romani$ and $\xvec^{\romanii}$.%
\footnote{%
  \cite{gumerov2013efficient} \textsc{Gumerov} \& \textsc{Duraiswami} (2013), Appendix B.1
}
In contrast to the earlier treatment of parametrization of filaments, the arclength shall not serve as the parameter, us rather choosing $t \in [-1,1]$.
Since the filament is a straight line, we have that $\xvec = \bkha + \sfrac{t}{2} (\xvec^{\romani} - \xvec^{\romanii})$.
It can then be shown that the velocity induced by a straight vortex filament is given by

\begin{equation}\label{eq:Biot-Savart_projection}
  \uvec(\bzhe) = \frac{\Gamma}{4\pi} \frac{\rvec^{\romani} \times \rvec^{\romanii}}{r^{\romani} r^{\romanii} + \rvec^{\romani} \cdot \rvec^{\romanii}} \left( \frac{1}{r^{\romani}} + \frac{1}{r^{\romanii}} \right),
\end{equation}

\noindent
where $\rvec^\romani = \bzhe - \xvec^\romani$, and so on.
Thus, the induced velocity at any point not on the filament can be calculated analytically.
As the point of evaluation approaches the vortex filament, the scalar product in equation \eqref{eq:Biot-Savart_projection} approaches $-r^\romani r^\romanii$, yielding infinite induced velocity, corresponding to the \textsc{Biot}--\textsc{Savart} kernel being singular.
A remedy for the implementation for such an expression is to add a small correction proportional to the difference in the denominator,%
\footnote{%
  \cite{kebbie-anthony2019fast} \textsc{Kebbie-Anthony} \emph{et al.} (2019), \S{II.C}
}
or implementing a smoothing kernel for the vorticity two-form.%
\footnote{%
  \cite{vatistas1991simpler} \textsc{Vatistas}, \textsc{Kozel} \& \textsc{Mih} (1991)
}
Opting for the former, one will write

\begin{equation}
  \uvec(\bzhe) = \frac{\Gamma}{4\pi} \frac{(\rvec^{\romani} \times \rvec^{\romanii}) (r^{\romani} + r^{\romanii})}{r^{\romani}r^{\romanii}(r^{\romani} r^{\romanii} + \rvec^{\romani} \cdot \rvec^{\romanii}) + {(\varepsilon \absl{\rvec^{\romanii} - \rvec^{\romani}})}^2},
\end{equation}

\noindent
where $\varepsilon$ is the smoothing parameter.
Figure \ref{fig:smoothing-comparison} illustrates the impact of the smoothing parameter, both in terms of its magnitude, and the distance of the point of evaluation from some vortex filament.
Ascending from some isoline on the horizontal axis, it is evident that placing the point of evaluation close the the filament will increase the relative error in the induced velocity.
In the same manner, increasing the smoothing parameter, perhaps better labeled the oft-called cutoff radius, increases the relative error overall.

\begin{figure}[h]
  \centering
  \includegraphics[scale = 1]{smoothing-comparison.pdf}
  \caption{The relative error in induced velocity for different smoothing parameters, evaluated at varying distances from some vortex filament in a wind turbine wake.}
  \label{fig:smoothing-comparison}
\end{figure}

Considering once more the mollified vorticity of equation \eqref{eq:mollified-vorticity}, we can investigate a mollifier which trades the compact support for regularity at the filament.
A collection of such possible mollifying functions have been discussed extensively in the literature.%
\footnote{%
  \cite{leonard1985computing} \textsc{Leonard} (1985), eq.17
}\textsuperscript{,}%
\footnote{%
  \cite{hoydonck2013computation} \textsc{van Hoydonck}, p.58
}
We shall employ the mollifier of \textsc{Rosenhead} and \textsc{Moore}, which is the first of such modifications to the kernel conceived,%
\footnote{%
  \cite{rosenhead1930spread} \textsc{Rosenhead} (1930)
}\textsuperscript{,}%
\footnote{%
  \cite{moore1972motion} \textsc{Moore} \& \textsc{Saffman} (1972)
}
the actual expression adapted from \textsc{Winckelmans} and \textsc{Leonard}.%
\footnote{%
  \cite{winckelmans1993contributions} \textsc{Winckelmans} \& \textsc{Leonard} (1993), eq.41; tab.II
}
It is imperative that we be aware of the dimensionality of the problem at hand.
For the literature treats two inherently different approaches in vortex mechanics in close proximity---one will find threedimensional distributions and twodimensional distributions, and it is the latter we shall employ.
Recall that the vorticity be a two-form supported along a filament, us considering the flux through the surface area of some vortex tube thereof.
The \textsc{Rosenhead}--\textsc{Moore} mollifier is then

\begin{equation}
  \varphi_{\sigma}(n^1, n^2) = \frac{\sigma^2}{\pi{(\absl{n}^2 + \sigma^2)}^2}, \qquad \int_{\reals^2} \frac{\sigma^2}{\pi{(\absl{n}^2 + \sigma^2)}^2} \,\dee n^1 \dee n^2 = 1.
\end{equation}

\noindent
The velocity that results from adding this mollifier can be derived analytically,%
\footnote{%
  \cite{hoydonck2013computation} \textsc{van Hoydonck}, Appendix F.1
}
and is what will be employed in the numerical machinery calculating the induced velocity from vortex lattices in the form of a wind turbine wake.
It is clear that the length of the vortex filament is equal to the difference in length from the point of evaluation to either endpoint of the filament.
That is, $\svec = \xvec^{\romani} - \xvec^{\romanii} = \rvec^{\romanii} - \rvec^{\romani}$, with $s$ denoting the magnitude of the vector, being the length of the filament.
For clarity in the expression of the velocity, we then denote the distance from the evaluation point to the filament, and its distances to the endpoints as follows.

\begin{equation}
  a = \frac{\absl{\rvec^{\romani} \times \rvec^{\romanii}}}{s}, \qquad b = \frac{\rvec^{\romani} \cdot \svec}{s}, \qquad c = \frac{\rvec^{\romanii} \cdot \svec}{s}
\end{equation}

\noindent
The mollified velocity is then

\begin{equation}
  \uvec_{\sigma} (\bzhe) = \frac{\Gamma (\rvec^{\romani} \times \rvec^{\romanii})}{4 \pi s (\sigma^2 + a^2)} \left( \frac{c}{\sqrt{\sigma^2 + a^2 + c^2}} - \frac{b}{\sqrt{\sigma^2 + a^2 + b^2}} \right).
\end{equation}
